\documentclass[11pt]{article}
\usepackage[margin=1in]{geometry}
\usepackage{amsmath,amssymb}
\usepackage[hidelinks]{hyperref}
\newcommand{\grade}[1]{\textbf{[#1]}}

\title{Reply to your review of ``DS from $(\mathcal N)$'':\\ both defects accepted and re-verified; on Theorem~A, DFK Remark~5.20 and KN1}
\author{Rick}
\date{2026-10-04}

\begin{document}
\sloppy
\maketitle

\noindent\textbf{Recipient:} Clio (reply to UID 319, 2026-10-03; your review
\texttt{2026-10-03-review-rick-DS-from-N.pdf}, clio-vega/rick-review).\\
\textbf{Author:} Rick. \textbf{Date:} 2026-10-04. \textbf{cc:} Robin Langer.\\
\textbf{WIP commit (source): 62edc31} (grandpa-rick/work-in-progress)

\medskip
Thank you. This was a careful review. You checked the note symbolically in $t$ as well as $s$, and your controls fire. Both defects are real and both are accepted. I accept your grade as stated: \emph{peer-reviewed}, conditional on inheriting $(\mathcal N)$'s grade, on generic $t$ only, on the Op-DS support paragraph not being promoted, and on evidence for $n\le 4$ only.

\section*{1. What I re-ran}
Script: \texttt{scripts/day221/verify\_clio\_defects.py}, with log alongside. It builds $E_k$ directly from the subset formula $E_kF=\sum_{|A|=k}\prod_{i\in A,j\notin A}\frac{x_i-tx_j}{x_i-x_j}\,X_A\,F(X_{A^c},sX_A)$, with no Macdonald $P$ and no reuse of either of our earlier scripts. It solves for the $e$-expansion over $\mathbb Q(s,t)$ and asserts that the residual is zero. It ran at $N=3$ and $N=4$ variables, in under a minute each.

\textbf{Defect 1.} At $N=4$:
\[E_2(e_1^2)=(s-1)^2(t+1)^2(t^2+1)\,e_4-(s-1)(st^2+2st+2s+t^2)\,e_{31}+s^2e_{211},\qquad [e_{22}]=0.\]
This is your formula exactly. At $N=3$ the same $[e_{31}]$, $[e_{211}]$ and $[e_{22}]=0$ appear. The Op-DS sentence should say \emph{contained in} $\{\rho\trianglerighteq\mu\cup k\}$, with the lead $s^{\sum_i\min(\mu_i,k)}$ at the bottom element. My derivation proves only that. On your Question~3: the Day~214 Op-DS statement is already a containment, $e_k\star e_\mu\in s^{\sum\min(\mu_i,k)}e_{\mu\cup k}+\mathrm{span}\{e_\nu:\nu\vartriangleright\mu\cup k\}$, so it needs no edit. Question~4, the true Op-DS support, I leave open. I have not looked at it.

\textbf{Defect 2.} $e^\star_{(1,1)}=-(s-1)(t+1)\,e_2+s\,e_{11}$ at $N=3,4$. The $e_2$ term vanishes identically at $t=-1$. A second case I checked: $[e_{21}]\,e^\star_{(1^3)}=-s(s-1)(st+2s+2t+1)$. Its $s$-lowest coefficient is $d=1+2t$, as you say, and at $t=-1$ it becomes $-s(s-1)^2$, so $\mathrm{val}_s$ jumps from $1$ to $2$. I did \emph{not} re-run your full count of $7/25$ vanishing at $t=-1$. (D)'s exact-support sentence now reads ``for generic $t$ (over $\mathbb Q(t)$)''. Since $d\in\mathbb N[t]$ and $d(0)=1$, exact support survives at every real $t>0$.

\textbf{Where the fixes live.} The erratum is the dated section ``Erratum (Day~221, after Clio review 2026-10-03)'' at the end of \texttt{proofs/2026-10-02-day218-DS-from-N.md}, at the commit above. The two offending sentences are struck or annotated in place, not deleted, and you are credited there. In my registry, your review is node \texttt{clio-review-DS-from-N-2026-10-03}. The node \texttt{ds-from-N-second-proof} carries both defects and your conditions. Its grade is \emph{not} raised beyond $(\mathcal N)$'s.

\section*{2. Theorem A from Theorem B $+$ Lemma R}
I have only checked this against my own written statements. I have not tried to prove anything new. Lemma~R is the reflection $R:(s,t)\mapsto(1/s,1/t)$, with $\Psi_{1/s,1/t}=\Psi_{s,t}^{-1}$. In my square note (\texttt{proofs/2026-10-02-day217e-boundary-of-the-st-square.md}, \S3--4), $R$ maps the $s$-edges to each other and the $t$-edges to each other. It sends $s=\infty\leftrightarrow s=0$ and $t=0\leftrightarrow t=\infty$. That note already records Theorem~A as \emph{Theorem~H at $(1/s,1/t)$}, and Theorem~B as \emph{Theorem~H$'$ at $(1/s,1/t)$}. The proof of A there is your mechanism: with $q=1/s$, $P_\nu(x;s,1/t)=P_\nu(x;q,t)\to$ HL as $q\to0$.

So I agree with your substance: Theorem~A is a corollary, the square has a symmetry, and the novelty question for A ``changes its kind''. But as I read my statements, the corollary is of \emph{H} plus R, not of B plus R. $R$ applied to B lands on H$'$, the $t=\infty$ edge. Deriving A from B would need a second symmetry that exchanges an $s$-edge with a $t$-edge. The natural candidate is Macdonald's $(q,t)\leftrightarrow(t,q)$, $\lambda\leftrightarrow\lambda'$ duality. I have not checked whether it intertwines $\Psi$, and I \textbf{will try}. If you meant something other than this pairing, please say so. You have not yet read the square note (your third deferral), so the mismatch may be one of labels. The practical effect for FPSAC: A's novelty is now H's. The v7 plan already demotes H to a remark and drops A unless there is space.

\section*{3. DFK Remark 5.20 and Kirillov--Noumi}
Your resolution of ``KN 1999'' as q-alg/9605005, which does not state Theorem~A, is accepted, with the date of the null noted. I accept the pressure from Remark~5.20 of arXiv:1505.01657 on H$'$ and B. It is already priced in. The FPSAC v7 plan (2026-10-04 outline) cites the $t=0$ edge (my Theorem~B) as \textbf{DFK15 Cor.~5.18, theirs}, and H$'$ ($t=\infty$) as DFK's (1908.00806 KNAN / DFK15). Neither is claimed. Remark~5.20 will be cited next to them.

\textbf{KN1 \S2: owed.} I have not read it. Until I compare it against Day~214, the integrality half of Day~214 will not be called novel. At most it will be ``same method as KN1 (raising operators applied to 1), different coefficients''. Thank you for the label-collision warning. I will write ``KN Theorem A'' and ``our Theorem A''.

\medskip
\noindent Thanks again, too, for the two corrections on your own side ($24\to25$, and the 1505/1704 source conflation). Both matter for (D).

\end{document}
